\clearpage
\fancyhead[L]{\small BODE. ASTRONOMISCHES JAHRBUCH.}
\hypertarget{native-568-lower}{}
\begin{center}\small Göttingische gelehrte Anzeigen. Stück 40. Seite 393..400. 1814 März 10.\\[3pt]\rule{.72\linewidth}{.3pt}\end{center}

\emph{Astronomisches Jahrbuch für das Jahr} 1816, \emph{nebst einer Sammlung der neuesten in die astronomischen Wissenschaften einschlagenden Abhandlungen, Beobachtungen und Nachrichten. Berechnet und herausgegeben von} J.\,E. \textsc{Bode} königl. Astronom und Mitglied der Akademie. 1813. 268 Seiten in Octav.

Das Jahr 1816 zeichnet sich durch eine grosse Sonnenfinsterniss aus (19. November), welche in einem Theile von Europa total mit Dauer sein wird, z.\,B. in Danzig, Warschau. Die Einrichtung des Kalenders ist in diesem Jahre ganz dieselbe, wie in dem vorhergehenden, nur sind jetzt für Sonnendurchmesser und Schiefe der Ekliptik die neuesten Bestimmungen zum Grunde gelegt, und bei den Verfinsterungen der Jupiterstrabanten, welche noch nach den \textsc{Wargentin}'schen Tafeln berechnet sind, ist der Gang der Unterschiede dieser Tafeln von den \textsc{Delambre}'schen im Allgemeinen beigefügt. Von der Menge und Reichhaltigkeit der \emph{Zusätze}, welche auch diesen Band begleiten, wurden wir um so angenehmer überrascht, da wir der Zeitumstände wegen kaum gehofft hatten, das Jahrbuch überhaupt zur gewöhnlichen Zeit erscheinen zu sehen. Den Anfang macht ein chronologisches Verzeichniss der berühmtesten (verstorbenen) Astronomen seit dem dreizehnten Jahrhundert, ihrer Verdienste, Schriften und Entdeckungen, in welchen man nicht leicht einen in der Geschichte der Astronomie nur irgend merkwürdigen Namen vermissen wird. — Astronomische Beobachtungen auf der königl. Sternwarte zu Kopenhagen in den Jahren 1811 und 1812 angestellt, von Herrn Staatsrath und Ritter \textsc{Bugge}. Oppositionen des Uranus 1811 und 1812, des Saturn 1811 und 1812, der Ceres 1811, der Vesta 1812, des Mars 1811, Beobachtungen des grossen Cometen von 1811, Sternbedeckungen und Verfinsterungen der Jupiterstrabanten. — Beobachtete Jupiterstrabanten-Verfinsterungen und Sternbedeckungen auf der Greenwicher Sternwarte in den Jahren 1809 und 1810 (aus der gedruckten Sammlung der \textsc{Maskelyne}'schen Beobachtungen gezogen). — Beobachtungen des veränderlichen Sterns $\eta$ im Antinous, und Tafeln zur Berechnung seines grössten Lichts von Hrn. Prof. \textsc{Wurm}, eine Fortsetzung des Aufsatzes im Jahrbuche für 1814. Wir wünschten, dass dieser verdiente Astronom uns auch bald mit einer ähnlichen Arbeit über den veränderlichen Stern im Wallfisch beschenken möchte. — Astronomische Beobachtungen auf der königl. Sternwarte zu Berlin angestellt, im Jahre 1812. — Beobachtungen der Jupiterstrabanten. — Verfinsterungen, Sternbedeckungen, Gegenscheine der Ceres, Pallas, Juno, Vesta, des Cometen von 1812 auf der Wiener Sternwarte angestellt von Hrn. \textsc{Triesnecker}. Meridianbeobachtungen der Ceres und Vesta, des Mars, Saturn und Jupiter im Jahre 1811 angestellt von Hrn. \textsc{Groombridge}. — Ueber die Bestimmung der Theilungsfehler eines Spiegelsextanten von Hrn. Dr. \textsc{Benzenberg}. Dass man sich auf die absoluten Grössen der mit diesen Werkzeugen gemessenen Winkel nicht unbedingt verlassen dürfe, ist sehr gegründet, wenn es gleich in ein\SourceLeaf{569}zelnen Fällen unentschieden bleiben mag, wie viel von dieser Ungewissheit auf Rechnung wirklicher Theilungsfehler, oder einer kleinen Excentricität oder einer kleinen Abweichung der beiden Flächen des grossen Spiegels von der parallelen Lage komme. Das von Hrn. B. vorgeschlagene Verfahren, die Fehler des Instruments in eine Tafel zu bringen, verträgt indess nur in sehr ebnen Gegenden eine Anwendung, wenn man nicht zugleich ein Instrument zum Messen kleiner Höhen hat; auch ist es wohl nicht überflüssig, dabei eine besondere Vorsicht zu empfehlen, damit alle Winkel genau aus Einem Punkt gemessen oder wenigstens darauf reducirt werden. — Verschiedene astronomische Beobachtungen, von Hrn. \textsc{Lee} aus London mitgetheilt. Beobachtungen des grossen Cometen von 1811 im Ostindischen Ocean am Bord eines Schiffes vom 17. Mai bis 15. Juni; ferner verschiedene astronomische Beobachtungen aus Port Jackson auf Neuholland und aus Calcutta. — Astronomische Beobachtungen im Jahre 1812 auf der Prager Sternwarte von den Hrn. \textsc{David} und \textsc{Bittner}. — Beobachtete Scheitelabstände der Sonne und Sternbedeckungen zu St. Gallen in der Schweiz, von Hrn. \textsc{von Scherer}, woraus Herr \textsc{Triesnecker} zugleich die geographische Lage dieses Orts ableitet. — Sichtbare Lichtveränderungen des Algol in den Jahren 1814, 1815 und 1816, vorausberechnet von Hrn. \textsc{Wurm}. — Beobachtungen des grossen Cometen von 1811 vom 6. September bis 19. October, Oppositionen des Saturn 1811 und Sternbedeckungen von Hrn. \textsc{Derflinger} in Kremsmünster. — Neue Refractionstafel, aus \textsc{Bradley}'s Beobachtungen der Circumpolarsterne abgeleitet, von Hrn. Prof. \textsc{Bessel}. Diese Tafel ist aus einer Abhandlung des Verf. im vierten Heft des Königsberger Archivs gezogen; allein es fehlt dabei eine Anweisung zum Gebrauch, um so mehr, da das Rechnungsbeispiel nicht ganz im Geist des Verf. ausgeführt ist. — Beobachtungen des grossen Cometen im Jahre 1811, nebst Bemerkungen über den Bau seiner verschiedenen Theile, von Herrn Dr. \textsc{Herschel}; und nachher Beobachtungen des zweiten Cometen vom Jahre 1811 nebst Bemerkungen über seinen Bau; von demselben Verf. sind Auszüge aus zwei Abhandlungen in den neuesten Bänden der \emph{Philosophical Transactions}, von denen wir hier einer umständlichen Anzeige um so mehr uns überheben, da die Originalabhandlungen bereits früher in diesen Blättern von einem andern Rec. angezeigt sind. Wir bemerken daher nur, dass, so schätzbar die Beobachtungen \textsc{Herschel}'s über diese Cometen sind, es uns jetzt noch viel zu früh scheint, Resultate über die Naturgeschichte dieser merkwürdigen Weltkörper feststellen zu wollen. Der Erfahrungen sind noch viel zu wenige, sie haben selbst bei weitem nicht den Grad von Zuverlässigkeit, wie astronomische Beobachtungen anderer Art, und unsrer Ueberzeugung nach werden erst die künftigen Jahrhunderte, die die jetzt beobachteten Cometen werden wiederkehren sehen, über das, was jetzt nur vage Hypothese sein kann, entscheiden können. Herr \textsc{Herschel} stellt unter andern die Vermuthung auf, dass der Comet von 1807, ehe er diesmal in sein Perihelium gelangte, früher sich einmal einem andern Fixsterne genähert habe und dadurch gleichsam zu einer frühern Reife gekommen sei als der von 1811, und glaubt, dass Cometen um andere Sonnen als die unsrige laufen mögen, werde durch den Umstand sehr wahrscheinlich, dass wir unter der grossen Anzahl der bis jetzt beobachteten erst von einem einzigen die Rückkehr mit Gewissheit kennen. Allein hiebei übersieht er offenbar, dass nach allen bisherigen Resultaten die Umlaufszeiten der Cometen in der Regel eher noch nach Jahrtausenden als nach Jahrhunderten gemessen werden müssen, und dass erst seit ein Paar Jahrhunderten die Cometen ordentlich beobachtet und seit nicht viel mehr als einem Jahrhundert berechnet werden; und was namentlich den Cometen von 1807 betrifft, so weiss man, dass dessen grösste Distanz von der Sonne, eben so wie die grösste Distanz aller andern, deren Bahn elliptisch hat berechnet werden können, nicht viel mehr als Nichts ist gegen die Distanz der nächsten Fixsterne. — Sternbedeckungen, Jupiterstrabanten-Verfinsterungen, Mondfinsterniss, Vesta, Uranus und Mars beobachtet auf der Sterwarte in Wilna in den Jahren 1811, 1812, 1813 von Hrn. \textsc{Joh. Sniadecky}. — Beobachtungen des grossen Cometen von 1811 auf der Sternwarte zu Palermo und Berech\SourceLeaf{570}nung der Elemente seiner Bahn von Hrn. \textsc{Piazzi} aus dessen gedruckter Abhandlung. — Aus einem Schreiben des Hrn. Prof. \textsc{David} in Prag; bezieht sich auf die abweichenden Resultate, welche die Hrn. \textsc{Vent}, \textsc{Bode} und \textsc{David} zu verschiedenen Zeiten für die Polhöhe der Böhmischen Riesenkuppe gefunden hatten, worüber natürlich uns hier kein Urtheil zustehen kann. — Nachtrag zu dem im Jahrbuche 1813 befindlichen Verzeichnisse sichtbarer Sternbedeckungen im Jahre 1813 von Hrn. \textsc{von Wisniewsky}. — Verzeichniss der Länge und Breite von neun der vornehmsten Fixsterne nach den neuesten Beobachtungen, aus dem \emph{Nautical Almanac} für 1815 entlehnt. — Geographische Lage verschiedener Oerter im mittlern Amerika, bestimmt von Hrn. \textsc{von Humboldt} (aus dessen \emph{Recueil d'observations astronomiques etc.}). — Projectionsmethode einer allgemeinen Himmelskarte. Sie besteht darin, die Abstände vom Nordpol in der Projection den Tangenten von $\frac{2}{5}$ der Abstände auf der Kugel proportional zu setzen, indess bekanntlich in der stereographischen Projection jene den Tangenten der halben Abstände auf der Kugel proportional sind. Nach dieser neuen Art hat der Verf. die Himmelskarte bei der neuesten Ausgabe seiner geschätzten Anleitung zur Kenntniss des gestirnten Himmels entworfen, deren früheren Ausgaben eine ähnliche Karte nach der stereographischen Projection beigegeben war. Er wählte diese Abänderung, damit die Grade nach dem Südpol zu nicht so unverhältnissmässig stark gegen die nördlichen anwachsen, aber freilich gehen dadurch manche Vortheile der stereographischen Entwerfungsart verloren; das Bild bleibt dem Original nicht mehr in \emph{den kleinsten Theilen} ähnlich, alle schief gegen den Aequator liegenden Kreise werden nicht mehr durch Kreise dargestellt, sondern durch Curven höherer Ordnung, die nur mühsam durch einzelne Punkte construirt werden können, und die Winkel, unter welchen sie sich schneiden, sind denen auf der Kugel nicht mehr gleich. — Geocentrischer Lauf der Pallas vom 1. August 1814 bis 1. Februar 1815, und der Vesta vom 18. October 1813 bis 9. Juli 1814. — Planetenbeobachtungen auf der Greenwicher Sternwarte 1809 und 1810 von Dr. \textsc{Maskelyne}. — Beobachtungen der beiden im Jahre 1813 erschienenen Cometen. Es sind hier blos vier Pariser Beobachtungen des ersten Cometen aufgeführt, und über den andern (von Hrn. Prof. \textsc{Harding} auf der hiesigen Sternwarte und zugleich von Hrn. \textsc{Pons} in Marseille entdeckten) nur einige allgemeine Nachrichten gegeben. — Die elliptischen Elemente der Planetenbahnen, für die ältern Planeten aus der dritten Ausgabe der \textsc{Laplace}\emph{'schen exposition du système du monde} (wovon 1812 eine vierte Ausgabe erschienen ist), die der neuern nach den neuesten (dem Herausgeber damals bekannten) Bestimmungen von \textsc{Gauss}; dieser Artikel würde jetzt bereits eine fast gänzliche Umarbeitung vertragen, da die Theorien der meisten Planeten seitdem neu bearbeitet sind. — Fernere Beobachtungen des grossen Cometen von 1811 nebst astronomischen Bemerkungen von Hrn. Prof. \textsc{Bessel}, eine neue Reduction der wichtigen \textsc{von Zach}'schen Beobachtungen dieses Cometen im ersten Zweige seiner Bahn. — Beobachtungen des Cometen von 1812 nebst Elementen seiner Bahn von \textsc{Bouvard}. — Methode zur Bestimmung der Abweichung eines Passageinstruments vom Meridian aus Beobachtungen der obern und untern Culminationen zweier nordlichen in der geraden Aufsteigung beinahe entgegengesetzten Sterne, nebst einem Verzeichniss solcher Sterne. Wir halten diese Methode für die allerzweckmässigste, besonders wenn der eine Stern der Polarstern selbst ist, und bemerken nur, dass in diesem Falle der andere Stern ohne Nachtheil auch von beinahe \emph{gleicher} Rectascension sein darf. — Planetenbeobachtungen auf der Pariser Sternwarte im Jahre 1809 von Hrn. \textsc{Bouvard} (aus der \emph{Connoissance des tems} für 1812). — Unter den am Schluss beigefügten verschiedenen kleinen astronomischen Notizen bezeichnen wir hier nur diejenigen, welche die vortrefflichen aus der \textsc{Reichenbach}'sche Werkstatt in München hervorgegangenen astronomischen Instrumenten betreffen, worunter ein 24fussiger Achromat mit 8 Zoll Oeffnung. — Noch sind zwei Nachträge zu diesem Bande zu bemerken, wovon der erstere die Originalbeobachtungen des grossen Cometen von 1811, von \textsc{Wisniewsky} im August 1812 angestellt, der andere, das \textsc{Pond}'sche Fixsternverzeichniss (obwohl nicht \SourceLeaf{571}das allerneueste) enthält. — Wir können nicht umhin, dem verdienstvollen Herausgeber für den unermüdeten Eifer zu danken, womit er so vielfache dem Astronomen wichtige Materialien zusammengetragen hat, unter denen auch dasjenige, was aus andern, dem bei weiten grössten Theile der Leser ohnehin jetzt meistens nicht zugänglichen, Schriften entlehnt ist, nicht anders als willkommen sein kann.
